\documentclass[12pt]{article}
\usepackage[utf8]{inputenc}
\usepackage[T1]{fontenc}
\usepackage[margin=1in]{geometry}
\usepackage{setspace}
\usepackage{xcolor}
\usepackage[colorlinks=true, urlcolor=blue!60!black, linkcolor=blue!60!black]{hyperref}
\setlength{\parindent}{2em}
\setlength{\parskip}{1.25em}
\renewcommand{\baselinestretch}{1.0}
\begin{document}
\begin{center}
{\LARGE \textit{McCulloch v. Maryland}, 17 U.S. (4 Wheat.) 316 (1819)\par}
\vspace{0.8em}
{\large Case Brief --- Four-Page Edition\par}
\end{center}
\vspace{1em}

\section{Memory Jogger}

The Necessary and Proper Clause gives Congress implied authority to use any appropriate, unprohibited means, including chartering a national bank, to execute its enumerated powers, and the Supremacy Clause forbids states to tax or impede federal instrumentalities (\href{https://supreme.justia.com/cases/federal/us/17/316/}{17 U.S. 316}).

\section{Detailed Case Facts}

After the First Bank of the United States expired and the War of 1812 disrupted federal finances, Congress chartered the Second Bank of the United States on April 10, 1816. In 1817 the Bank opened a Baltimore branch without any authorization from Maryland. Resenting its competition with state banks and its contraction of credit in the downturn of 1818, Maryland enacted a tax on all banks ``not chartered by the legislature,'' a category that in practice contained only the Baltimore branch. Unchartered banks could issue notes only on stamped paper bought at steep per-note fees, unless they paid \$15,000 a year in advance; violating officers forfeited \$500 per offense.

James William McCulloch, the branch's cashier, issued unstamped notes without paying the fee. John James sued him as a \textit{qui tam} informer for himself and the State.

\section{Procedural History}

On an agreed statement of facts, the Baltimore County Court entered judgment against McCulloch, and the Maryland Court of Appeals affirmed. On writ of error, after nine days of argument featuring Daniel Webster, William Pinkney, and William Wirt for McCulloch and Luther Martin for Maryland, the Supreme Court reversed on March 6, 1819, and ordered judgment for McCulloch.

\section{Judicial Votes}

Reversed, \textbf{unanimously}.
\begin{itemize}
\item \textbf{Opinion of the Court:} Chief Justice Marshall, joined by Justices Washington, Johnson, Livingston, Duvall, and Story. Standard databases list Justice Todd as joining a 7--0 decision; some historical works report him absent, making the vote 6--0.
\item \textbf{Concurring or dissenting:} None.
\end{itemize}

\section{Holding}

(1) Congress may choose any appropriate, unprohibited means plainly adapted to executing its enumerated powers. ``Necessary'' does not mean strictly indispensable but includes means that are convenient, useful, or conducive to those ends, and the degree of necessity is for Congress, so long as it does not act under a ``pretext.'' Congress could therefore charter the Bank and its branches. (2) States have no power, by taxation or otherwise, to retard, impede, burden, or control the operations of constitutional federal laws, so Maryland's tax is void. The pretext limitation and the carve-out permitting nondiscriminatory taxes on the Bank's real property and on Marylanders' shares were dicta.

\section{Analysis of Opinions}

\textbf{Part I: the power to charter the Bank.} Marshall began with practice: the First Congress debated and chartered the First Bank, successive governments recognized it, and even former skeptics were persuaded by the crises after its expiration. He rejected Maryland's compact theory: the Constitution was ratified by conventions of the people, making the national government ``emphatically and truly, a Government of the people,'' supreme within its sphere.

Although ``bank'' and ``incorporation'' are not enumerated, the Tenth Amendment, unlike the Articles of Confederation, omits the word ``expressly,'' leaving room for implied powers. A constitution cannot have the prolixity of a legal code; it marks great outlines, and the means follow from the great powers of the sword and the purse. Incorporation is not a great substantive and independent power but an ordinary means.

Marshall then rejected Maryland's narrow reading of ``necessary'' on four grounds. In ordinary usage it often means convenient or useful, and Article I, Section 10's ``absolutely necessary'' shows the Framers knew how to demand strict necessity. A strict rule would be unworkable, invalidating even federal punishment of mail robbery. The Clause appears among Congress's grants of power, not its limitations. And once a means is appropriate and not pretextual, its necessity is for Congress to judge.

\textbf{Part II: Maryland's tax.} States' concurrent taxing power reaches only subjects within their sovereignty, not instruments created by the people of all the states. Because the power to tax involves the power to destroy, a single state taxing a federal instrument would let a part control the whole without the check of representation, since Maryland's legislators answer only to Marylanders. Refusing to judge how heavy a tax is too heavy, Marshall drew a bright line against state taxation of federal operations.

\section{Examples: Future Applications}

\begin{enumerate}
\item \textbf{Same-Side (Valid): Civil Commitment.} Congress authorizes civil commitment of sexually dangerous federal prisoners after their sentences end. Running federal prisons is a means of executing enumerated powers, as \textit{United States v.\ Comstock} held.
\item \textbf{Same-Side (Invalid State Tax): A Postal Surcharge.} A state imposes a \$50 surcharge on Postal Service deliveries while exempting private couriers. It directly burdens a federal instrumentality.
\item \textbf{Opposite-Side (Invalid Federal Law): An Insurance Mandate.} Under Chief Justice Roberts's analysis in \textit{NFIB v.\ Sebelius}, a mandate to buy insurance may be necessary to a regulatory scheme but is not proper, because compelling commerce is a great substantive and independent power.
\item \textbf{Opposite-Side (Valid State Tax): General Taxes.} A nondiscriminatory property tax on a federally chartered bank's building, or a general income tax on federal employees, falls within \textit{McCulloch}'s carve-out because it burdens the state's own constituents too.
\item \textbf{Fence-Sitter: A Chain of Implied Powers.} Congress makes it a felony for a subgrantee's employee to falsify any local business record, even where no federal funds are misused. Proponents cite \textit{Comstock} and \textit{Sabri}; challengers say stacked inferences stretch ``plainly adapted'' and invade state criminal jurisdiction.
\end{enumerate}

\section{Critique}

\textbf{Originalist and textualist views.} Spencer Roane immediately attacked the opinion for dissolving enumerated limits, prompting Marshall's pseudonymous defense as ``A Friend of the Constitution.'' Gary Lawson, Geoffrey Miller, Robert Natelson, and Guy Seidman argue from eighteenth-century agency law that ``necessary'' required a tighter fit and ``proper'' imposed independent limits; William Baude stresses Madison's line between incidental and great substantive powers, revived in \textit{NFIB}. Akhil Reed Amar and Jack Balkin defend Marshall's intratextual reading.

\textbf{Progressive and structural views.} John Mikhail and David Schwartz argue that \textit{McCulloch} conceded too much to strict enumerationism, with Mikhail reading the Sweeping Clause to reach powers vested in the government as a whole. The case became the foundation for the administrative state and New Deal regulation, and its representation rationale for tax immunity remains one of the most influential structural arguments in constitutional law.

\section{Key Quotations}

\begin{quote}
``Let the end be legitimate, let it be within the scope of the Constitution, and all means which are appropriate, which are plainly adapted to that end, which are not prohibited, but consist with the letter and spirit of the Constitution, are Constitutional.''
\end{quote}
\noindent --- Chief Justice Marshall, Opinion of the Court, 17 U.S. (4 Wheat.) at 421.

\begin{quote}
``That the power to tax involves the power to destroy; that the power to destroy may defeat and render useless the power to create; that there is a plain repugnance in conferring on one Government a power to control the constitutional measures of another, which other, with respect to those very measures, is declared to be supreme over that which exerts the control, are propositions not to be denied.''
\end{quote}
\noindent --- Chief Justice Marshall, Opinion of the Court, 17 U.S. (4 Wheat.) at 431.

\begin{quote}
``Its nature, therefore, requires that only its great outlines should be marked, its important objects designated, and the minor ingredients which compose those objects be deduced from the nature of the objects themselves. \ldots{} In considering this question, then, we must never forget that it is a Constitution we are expounding.''
\end{quote}
\noindent --- Chief Justice Marshall, Opinion of the Court, 17 U.S. (4 Wheat.) at 407.

\end{document}
